\documentclass[10pt,letterpaper]{article}
\usepackage{cornell-notes}
\usepackage{mathtools}

\title{Chapter 16: Classification and Inference}
\author{}
\date{}

\setDocTitle{Chapter 16: Classification and Inference}
\setDocSubtitle{Numerical Methods}
\setDocVersion{v1.0}
\setDocStatus{Study Companion}
\setDocOwner{Numerical Methods Cornell Notes}
\setDocAuthor{}
\setDocDate{\today}
\setDocSummary{Cornell-format study and review notes for Chapter 16: Classification and Inference.}

\setCornellCollection{Numerical Methods Cornell Notes}
\setCornellUnitType{Chapter}
\setCornellUnitNumber{16}
\setCornellUnitTitle{Classification and Inference}
\setCornellCompanionLabel{Study and review companion}

\hypersetup{pdftitle={Chapter 16: Classification and Inference},pdfsubject={Cornell notes},pdfauthor={}}

\begin{document}
\maketitle

\section*{Chapter Overview}
This chapter organizes the mathematical and computational ideas behind classification and inference. The notes preserve the numbered section sequence while emphasizing method selection, explicit assumptions, numerical reliability, cost, and verification.

\section*{Learning Objectives}
Explain the governing problem class; compare Gaussian Mixture Models and k-Means Clustering, Viterbi Decoding, Markov Models and Hidden Markov Modeling, and Hierarchical Clustering by Phylogenetic Trees; design defensible stopping and verification checks; distinguish problem conditioning from algorithmic stability.

\section*{Prerequisites}
Probability, optimization, distances, dynamic programming, and linear algebra.

\section*{Operational Workflow}
Define labels or latent states and a loss; choose a generative, margin, clustering, or sequence model; fit without leakage; tune on validation data; assess calibration and out-of-sample error.

\subsection*{Formula and notation anchor}
\[
\widehat{c}(x)=\arg\max_c P(c\mid x),\qquad P(c\mid x)\propto p(x\mid c)P(c)
\]
A maximum-posterior decision is optimal only for the encoded loss; unequal costs require a cost-sensitive rule.

\begin{warnbox}{Numerical warning}
Label switching, local optima, poor feature scaling, invalid distance metrics, underflow in probability products, overfitting kernels, class imbalance, and treating clusters as ground truth.
\end{warnbox}

\section{16.0 Introduction}
\begin{cornell}
\noterow{What is the central idea?}{Classification and inference map observations to latent groups, paths, or labels while separating model fit from out-of-sample performance.}
\noterow{How should it be used?}{For Introduction, begin from explicit inputs, outputs, preconditions, and representation choices.  Define labels or latent states and a loss; choose a generative, margin, clustering, or sequence model; fit without leakage; tune on validation data; assess calibration and out-of-sample error.}
\noterow{Cost and reliability}{Analyze arithmetic work, storage, data movement, and convergence separately.  Relevant failure pressures include label switching, local optima, poor feature scaling, invalid distance metrics, underflow in probability products, overfitting kernels, class imbalance, and treating clusters as ground truth.  Use scaled tolerances and report both success criteria and diagnostic evidence.}
\noterow{Implementation checklist}{\begin{itemize}
\item Validate dimensions, domains, ordering, and structural assumptions.
\item Guard small divisors, invalid states, overflow, underflow, and nonfinite values.
\item Make mutation, workspace, indexing, and termination behavior explicit.
\item Test a simple exact case, a difficult valid case, a boundary case, and an expected failure.
\end{itemize}}
\end{cornell}
\begin{summarybox}{Section summary}
Classification and inference map observations to latent groups, paths, or labels while separating model fit from out-of-sample performance.  Select this technique only when its assumptions match the data, and accept a result only after an independent residual, error estimate, invariant, or refinement check.
\end{summarybox}

\section{16.1 Gaussian Mixture Models and k-Means Clustering}
\begin{cornell}
\noterow{What is the central idea?}{Mixture models estimate soft probabilistic assignments; k-means is a hard-assignment limit focused on within-cluster squared distance.}
\noterow{How should it be used?}{For Gaussian Mixture Models and k-Means Clustering, begin from explicit inputs, outputs, preconditions, and representation choices.  Define labels or latent states and a loss; choose a generative, margin, clustering, or sequence model; fit without leakage; tune on validation data; assess calibration and out-of-sample error.}
\noterow{Quantitative anchor}{\[
\min_{\mu,c}\sum_i\|x_i-\mu_{c_i}\|_2^2
\]
State the dimensions, scaling, and validity assumptions before using this relation.  Check the resulting residual or invariant rather than trusting an iteration count alone.}
\noterow{Cost and reliability}{Analyze arithmetic work, storage, data movement, and convergence separately.  Relevant failure pressures include label switching, local optima, poor feature scaling, invalid distance metrics, underflow in probability products, overfitting kernels, class imbalance, and treating clusters as ground truth.  Use scaled tolerances and report both success criteria and diagnostic evidence.}
\noterow{Implementation checklist}{\begin{itemize}
\item Validate dimensions, domains, ordering, and structural assumptions.
\item Guard small divisors, invalid states, overflow, underflow, and nonfinite values.
\item Make mutation, workspace, indexing, and termination behavior explicit.
\item Test a simple exact case, a difficult valid case, a boundary case, and an expected failure.
\end{itemize}}
\end{cornell}
\begin{summarybox}{Section summary}
Mixture models estimate soft probabilistic assignments; k-means is a hard-assignment limit focused on within-cluster squared distance.  Select this technique only when its assumptions match the data, and accept a result only after an independent residual, error estimate, invariant, or refinement check.
\end{summarybox}

\section{16.2 Viterbi Decoding}
\begin{cornell}
\noterow{What is the central idea?}{Dynamic programming finds the most probable hidden-state path by retaining the best predecessor for every state and time.}
\noterow{How should it be used?}{For Viterbi Decoding, begin from explicit inputs, outputs, preconditions, and representation choices.  Define labels or latent states and a loss; choose a generative, margin, clustering, or sequence model; fit without leakage; tune on validation data; assess calibration and out-of-sample error.}
\noterow{Quantitative anchor}{\[
\delta_t(j)=p(y_t\mid j)\max_i\{\delta_{t-1}(i)a_{ij}\}
\]
State the dimensions, scaling, and validity assumptions before using this relation.  Check the resulting residual or invariant rather than trusting an iteration count alone.}
\noterow{Cost and reliability}{Analyze arithmetic work, storage, data movement, and convergence separately.  Relevant failure pressures include label switching, local optima, poor feature scaling, invalid distance metrics, underflow in probability products, overfitting kernels, class imbalance, and treating clusters as ground truth.  Use scaled tolerances and report both success criteria and diagnostic evidence.}
\noterow{Implementation checklist}{\begin{itemize}
\item Validate dimensions, domains, ordering, and structural assumptions.
\item Guard small divisors, invalid states, overflow, underflow, and nonfinite values.
\item Make mutation, workspace, indexing, and termination behavior explicit.
\item Test a simple exact case, a difficult valid case, a boundary case, and an expected failure.
\end{itemize}}
\end{cornell}
\begin{summarybox}{Section summary}
Dynamic programming finds the most probable hidden-state path by retaining the best predecessor for every state and time.  Select this technique only when its assumptions match the data, and accept a result only after an independent residual, error estimate, invariant, or refinement check.
\end{summarybox}

\section{16.3 Markov Models and Hidden Markov Modeling}
\begin{cornell}
\noterow{What is the central idea?}{A hidden Markov model combines transition dynamics with observation likelihoods, enabling filtering, smoothing, decoding, and parameter learning.}
\noterow{How should it be used?}{For Markov Models and Hidden Markov Modeling, begin from explicit inputs, outputs, preconditions, and representation choices.  Define labels or latent states and a loss; choose a generative, margin, clustering, or sequence model; fit without leakage; tune on validation data; assess calibration and out-of-sample error.}
\noterow{Quantitative lens}{Identify the governing equation, recurrence, probability law, or geometric predicate for Markov Models and Hidden Markov Modeling.  Define every variable and scale; then derive a residual, error estimate, or invariant that can be checked independently.}
\noterow{Cost and reliability}{Analyze arithmetic work, storage, data movement, and convergence separately.  Relevant failure pressures include label switching, local optima, poor feature scaling, invalid distance metrics, underflow in probability products, overfitting kernels, class imbalance, and treating clusters as ground truth.  Use scaled tolerances and report both success criteria and diagnostic evidence.}
\noterow{Implementation checklist}{\begin{itemize}
\item Validate dimensions, domains, ordering, and structural assumptions.
\item Guard small divisors, invalid states, overflow, underflow, and nonfinite values.
\item Make mutation, workspace, indexing, and termination behavior explicit.
\item Test a simple exact case, a difficult valid case, a boundary case, and an expected failure.
\end{itemize}}
\end{cornell}
\begin{summarybox}{Section summary}
A hidden Markov model combines transition dynamics with observation likelihoods, enabling filtering, smoothing, decoding, and parameter learning.  Select this technique only when its assumptions match the data, and accept a result only after an independent residual, error estimate, invariant, or refinement check.
\end{summarybox}

\section{16.4 Hierarchical Clustering by Phylogenetic Trees}
\begin{cornell}
\noterow{What is the central idea?}{Hierarchical clustering repeatedly joins or divides groups under a linkage rule, producing a tree whose interpretation depends on the distance model.}
\noterow{How should it be used?}{For Hierarchical Clustering by Phylogenetic Trees, begin from explicit inputs, outputs, preconditions, and representation choices.  Define labels or latent states and a loss; choose a generative, margin, clustering, or sequence model; fit without leakage; tune on validation data; assess calibration and out-of-sample error.}
\noterow{Quantitative lens}{Identify the governing equation, recurrence, probability law, or geometric predicate for Hierarchical Clustering by Phylogenetic Trees.  Define every variable and scale; then derive a residual, error estimate, or invariant that can be checked independently.}
\noterow{Cost and reliability}{Analyze arithmetic work, storage, data movement, and convergence separately.  Relevant failure pressures include label switching, local optima, poor feature scaling, invalid distance metrics, underflow in probability products, overfitting kernels, class imbalance, and treating clusters as ground truth.  Use scaled tolerances and report both success criteria and diagnostic evidence.}
\noterow{Implementation checklist}{\begin{itemize}
\item Validate dimensions, domains, ordering, and structural assumptions.
\item Guard small divisors, invalid states, overflow, underflow, and nonfinite values.
\item Make mutation, workspace, indexing, and termination behavior explicit.
\item Test a simple exact case, a difficult valid case, a boundary case, and an expected failure.
\end{itemize}}
\end{cornell}
\begin{summarybox}{Section summary}
Hierarchical clustering repeatedly joins or divides groups under a linkage rule, producing a tree whose interpretation depends on the distance model.  Select this technique only when its assumptions match the data, and accept a result only after an independent residual, error estimate, invariant, or refinement check.
\end{summarybox}

\section{16.5 Support Vector Machines}
\begin{cornell}
\noterow{What is the central idea?}{An SVM selects a maximum-margin separator, using kernels for nonlinear boundaries and regularization for overlap or noise.}
\noterow{How should it be used?}{For Support Vector Machines, begin from explicit inputs, outputs, preconditions, and representation choices.  Define labels or latent states and a loss; choose a generative, margin, clustering, or sequence model; fit without leakage; tune on validation data; assess calibration and out-of-sample error.}
\noterow{Quantitative anchor}{\[
\min_{w,b}\frac12\|w\|^2\quad\text{s.t.}\quad y_i(w^Tx_i+b)\ge1
\]
State the dimensions, scaling, and validity assumptions before using this relation.  Check the resulting residual or invariant rather than trusting an iteration count alone.}
\noterow{Cost and reliability}{Analyze arithmetic work, storage, data movement, and convergence separately.  Relevant failure pressures include label switching, local optima, poor feature scaling, invalid distance metrics, underflow in probability products, overfitting kernels, class imbalance, and treating clusters as ground truth.  Use scaled tolerances and report both success criteria and diagnostic evidence.}
\noterow{Implementation checklist}{\begin{itemize}
\item Validate dimensions, domains, ordering, and structural assumptions.
\item Guard small divisors, invalid states, overflow, underflow, and nonfinite values.
\item Make mutation, workspace, indexing, and termination behavior explicit.
\item Test a simple exact case, a difficult valid case, a boundary case, and an expected failure.
\end{itemize}}
\end{cornell}
\begin{summarybox}{Section summary}
An SVM selects a maximum-margin separator, using kernels for nonlinear boundaries and regularization for overlap or noise.  Select this technique only when its assumptions match the data, and accept a result only after an independent residual, error estimate, invariant, or refinement check.
\end{summarybox}

\section*{Method comparison}
\begin{center}
\small
\begin{tabularx}{\linewidth}{>{\bfseries\raggedright\arraybackslash}p{0.22\linewidth}XX}
\toprule
Method or lens & Best use & Main caution \\
\midrule
k-means & Compact spherical clusters & Fast; scale-sensitive and hard assignments \\
Hidden Markov model & Sequential latent states & Transition and emission assumptions \\
Support vector machine & High-dimensional classification & Margin-based; kernel and regularization tuning \\
\bottomrule
\end{tabularx}
\end{center}

\section*{Worked example}
\begin{exambox}{Independent worked example}
With class means $0$ and $4$, equal priors, and equal unit variances, the one-dimensional Gaussian decision boundary is $x=2$.  A point $x=3$ is assigned to the second class.  Unequal priors, variances, or misclassification costs move the boundary, so the midpoint is not universal.
\end{exambox}

\section*{Chapter synthesis}
\begin{summarybox}{Chapter synthesis}
The chapter's methods should be viewed as a decision system rather than a list of routines.  Begin with the mathematical problem and its conditioning, exploit structure, select an algorithm with compatible stability and cost, and verify the computed answer with evidence that is independent of the internal iteration.  The recurring implementation discipline is: define labels or latent states and a loss; choose a generative, margin, clustering, or sequence model; fit without leakage; tune on validation data; assess calibration and out-of-sample error.
\end{summarybox}

\subsection*{Common mistakes and numerical hazards}
\begin{warnbox}{Common mistakes and numerical hazards}
Label switching, local optima, poor feature scaling, invalid distance metrics, underflow in probability products, overfitting kernels, class imbalance, and treating clusters as ground truth.  A second recurring mistake is reporting only a computed value: always retain tolerances, iteration counts, residuals, refinement behavior, and relevant structural checks.
\end{warnbox}

\subsection*{Retrieval practice}
\textbf{Questions}
\begin{enumerate}
\item What problem does Gaussian Mixture Models and k-Means Clustering address, and which assumption is easiest to violate?
\item What problem does Viterbi Decoding address, and which assumption is easiest to violate?
\item What problem does Markov Models and Hidden Markov Modeling address, and which assumption is easiest to violate?
\item What problem does Hierarchical Clustering by Phylogenetic Trees address, and which assumption is easiest to violate?
\item What problem does Support Vector Machines address, and which assumption is easiest to violate?
\item How do conditioning, stability, and the final residual play different roles in this chapter?
\end{enumerate}

\textbf{Brief answers}
\begin{enumerate}
\item Mixture models estimate soft probabilistic assignments; k-means is a hard-assignment limit focused on within-cluster squared distance. Recheck the section's domain, structure, scaling, and failure diagnostics.
\item Dynamic programming finds the most probable hidden-state path by retaining the best predecessor for every state and time. Recheck the section's domain, structure, scaling, and failure diagnostics.
\item A hidden Markov model combines transition dynamics with observation likelihoods, enabling filtering, smoothing, decoding, and parameter learning. Recheck the section's domain, structure, scaling, and failure diagnostics.
\item Hierarchical clustering repeatedly joins or divides groups under a linkage rule, producing a tree whose interpretation depends on the distance model. Recheck the section's domain, structure, scaling, and failure diagnostics.
\item An SVM selects a maximum-margin separator, using kernels for nonlinear boundaries and regularization for overlap or noise. Recheck the section's domain, structure, scaling, and failure diagnostics.
\item Conditioning describes sensitivity of the mathematical problem, stability describes error propagation by the algorithm, and the residual measures satisfaction of the computed equations; none is a universal substitute for the others.
\end{enumerate}

\subsection*{Mastery checklist}
\begin{itemize}
\item[$\square$] I can explain and select Gaussian Mixture Models and k-Means Clustering without relying on a memorized code listing.
\item[$\square$] I can explain and select Viterbi Decoding without relying on a memorized code listing.
\item[$\square$] I can explain and select Markov Models and Hidden Markov Modeling without relying on a memorized code listing.
\item[$\square$] I can explain and select Hierarchical Clustering by Phylogenetic Trees without relying on a memorized code listing.
\item[$\square$] I can explain and select Support Vector Machines without relying on a memorized code listing.
\item[$\square$] I can state a scaled stopping rule and an independent verification check.
\item[$\square$] I can identify at least one conditioning risk and one algorithmic stability risk.
\item[$\square$] I can estimate time, storage, and data-movement costs at the appropriate level.
\end{itemize}

\end{document}

